% !PRINTED 760
\weberpage{760}

With this it is proved:

\gesperrt{The number $\pi$ is a transcendental number. The ``quadrature of the circle'' cannot be solved by a geometric construction in which only algebraic curves and surfaces are employed.}

\begin{center}
\S.\,207.

\gesperrt{The general theorem of Lindemann on the exponential function.}
\end{center}

The transcendence of the numbers $e$ and $\pi$, which is hereby proved, is contained as a special case in a very general theorem on the exponential function that \gesperrt{Lindemann} announced in the memoir mentioned above, and of which a fully worked-out proof is contained in the cited memoir of \gesperrt{Weierstrass}. We will prove this theorem here in conclusion, using the same means that we have employed for the two special cases. The theorem runs as follows:

I. \gesperrt{There exists no equation of the form} \weq{1}{C_1e^{z_1} + C_2e^{z_2} + \cdots + C_me^{z_m} = 0,} \gesperrt{in which the coefficients $C_1,C_2,\ldots,C_m$ are algebraic numbers and the exponents $z_1,z_2,\ldots,z_m$ are algebraic numbers distinct from one another, unless all the coefficients $C_1,C_2,\ldots,C_m$ are equal to zero.}

In order to prove it, we first derive a lemma:

Let
\[
\alpha_1,\alpha_2,\ldots,\alpha_r
\]
be arbitrary real or imaginary quantities, \gesperrt{distinct however from one another}, and let
\[
A_1,A_2,\ldots,A_r
\]
be likewise arbitrary quantities that \gesperrt{do not all vanish}.

The same shall hold of the two sequences of quantities
\[
\begin{gathered}
\beta_1,\beta_2,\ldots,\beta_s \\
B_1,B_2,\ldots,B_s
\end{gathered}
\]
We denote by
\[
\gamma_1,\gamma_2,\ldots,\gamma_t
\]
%CONT>>
